% Einheit 3 von 4 – Einheit 3 (bank/lineare-gleichungen/e3.jsonl, zusammenbau v0.1)
%% TODO Prüfungswort Einheit 3: nicht in der Marken-Zeile
%% TODO Zeitmarke Einheit 3: Marken-Zeile nicht lesbar
%% TODO Zweigzeile Teil 1: was hier gelernt wird (Satz in Schülersprache aus dem Einheitstitel) – Einheit 3
%% TODO Einheitentitel 3 nicht in der Mappe lesbar
\einheitenkopf[e3]{Einheit 3 von 4 · Einheit 3}
\setcounter{aufgabe}{15}

%% TODO Titel: Ich-kann-Satz fehlt in der Bank (Platzhalter: Kettenname) – Nr. 16
%% TODO Anweisung: ein Satz über der ersten Teilaufgabe fehlt in der Bank – Nr. 16
\begin{aufgabe}{x beidseitig}
%% TODO \rechenplatz{n}: Schrittzahl des Grundfalls fehlt in der Bank (nur bei fester Schreibform, 2.3 b)
\begin{teile}
\teil $7x + 2 = 3x + 14$ – wo steht weniger $x$? \leerfeld
\end{teile}
\begin{gleichungsraster}[2]
\gl{5x = 2x + 12} &
\gl{6x + 5 = 2x + 25} \\
\gl{4x + 7 + x = 2x + 22} &
\gl{3(x + 2) = x + 14} \\
\gl{\frac{x}{2} + 4 = x - 1} &
\gl{2(x + 3) = 2x + 6} \\
\gl{\text{$5(x - 2) = 3x + 4$ – löse und mache die Probe.}} & \\
\end{gleichungsraster}
\end{aufgabe}

%% TODO Titel: Ich-kann-Satz fehlt in der Bank (Platzhalter: Kettenname) – Nr. 17
%% TODO Anweisung: ein Satz über der ersten Teilaufgabe fehlt in der Bank – Nr. 17
\begin{aufgabe}{Dezimalzahlen}
\begin{gleichungsraster}[2]
\gl{0{,}5x + 1{,}2 = 3{,}7} & \\
\end{gleichungsraster}
\end{aufgabe}

%% TODO Titel: Ich-kann-Satz fehlt in der Bank (Platzhalter: Kettenname) – Nr. 18
%% TODO Anweisung: ein Satz über der ersten Teilaufgabe fehlt in der Bank – Nr. 18
\begin{aufgabe}{x beidseitig – Fehler finden, Begründen}
\begin{teile}
\teil Nina rechnet: \rechnung{2x + 5 + 4x &= 23 \\ 11x &= 23} Finde den Fehler.
\teil Ole rechnet: \rechnung{15 - (x - 3) &= 2x \\ 15 - x - 3 &= 2x \\ 12 &= 3x \\ 4 &= x} Finde den Fehler.
\teil Paul rechnet: \rechnung{3(x + 1) &= 3x + 3 \\ 3x + 3 &= 3x + 3 &&\mid -3x \\ 3 &= 3} Er schreibt: „Die Lösung ist 0." Finde den Fehler.
\teil $x + 4 = x + 9$ – gibt es eine Lösung? Begründe ohne Rechnung.
\end{teile}
\end{aufgabe}
